\section{Use the chapter package for construction and submission}
The package is \texttt{035-h263}. Its native reading contains all 120 construction
steps, and its twelve working folders retain each earlier exercise. Begin with
SAD, then follow search, prediction, six-block coordination, transform,
quantization, block reconstruction, bit output, coefficient syntax, macroblock
syntax, picture construction and clip I/O. Keep your own working files: later
composition helpers carry your completed instructions and required data forward.
A new shell is an exercise boundary, not a reason to discard a working dependency.

Inside the extracted package, examples that previously began with
\path{drafts/} begin with \path{construction/}. Package-native reading fragments
use those exported paths. The book's directly edited include files remain the
authoritative prose; package metadata does not regenerate them. Use the matching
stage helper and preserve its data boundaries when preparing another observation.
Custom source organization can satisfy the same interfaces without using those
helpers.

The small \path{starter/interfaces.asm} exposes the thirteen callable entries
without an implementation. \path{starter/main.asm} supplies the final clip caller
and its unfinished high-level dependency. During guided work, begin with the
matching smaller stage instead. Complete the public contracts in either mode;
the grader does not require private labels, scratch registers, exact instruction
counts or one valid entropy representation.

Run one completed domain while developing it:
\begin{verbatim}
python3 grade.py my-working.asm --domain sad \
    --java /path/to/java --cli /path/to/mars-cli.jar \
    --report sad-feedback.json
\end{verbatim}
Omit \texttt{--domain} for all twelve completed-interface domains after composing
one complete source. The runner supplies separate high-address grading storage.
Reserve the \texttt{h263Check\_}, \texttt{h263Motion\_},
\texttt{h263Prediction\_}, \texttt{h263Macro\_}, \texttt{h263Transform\_},
\texttt{h263Quant\_}, \texttt{h263Block\_}, \texttt{h263Writer\_},
\texttt{h263Coeff\_}, \texttt{h263Syntax\_}, \texttt{h263Picture\_} and
\texttt{h263Clip\_} prefixes for feedback. Do not paste the grading harness into
your program. Keep ordinary student data out of its reserved wide-address arena.
The packaged constants and ordinary callers fit the smaller normal data region.

The finite public suite contains 1196 cases across the twelve domains. Passing
the final clip driver alone is insufficient: an implementation can encode the
worked clip while mishandling negative residuals, unusual strides or a domain
boundary. The instructor distribution adds private boundary and repeated-call
checks. Run the complete inverse-transform accuracy gate on submitted transforms
as well as finite block feedback. Its six Annex A cohorts and exact-zero calls
check error limits rather than requiring one private implementation's rounding.

The package records the minimum MARS grader revision for independently targeted
output failures. Use a matching CLI archive and plugin build; an older assignment
runtime cannot accept an option it does not implement. Reports and prepared run
directories require new paths so earlier work remains inspectable.

\section{Submit evidence for the program you actually ran}
Submit the complete assembly source with all thirteen entries, the raw-input
recipe and its parameters, public feedback, inverse-accuracy results, and a
successful independent stream report. Retain the generated input metadata and
source/tool hashes. Include the encoded stream and reconstruction trace for the
named eight-picture fixture, together with the guest-work measurement boundary.
Do not substitute output from a host encoder for output from the submitted MIPS
program.

Explain one motion prediction, one quantization/reconstruction trace and one
failure path using actual values from your run. State when a buffer becomes an
accepted reference and why an output prefix does not prove overall success.
Compare equivalent correct implementations before making a performance claim.
A smaller file, lower guest count and shorter wall time are different observations.

Guided and free-writing work share these semantic requirements. Instructors may
control prerequisite checks, quiz progression, hint availability and mode switches;
remote delivery and automatic enforcement remain a separate plugin/service design.
The instructor materials contain explanations and progressive hints, while the
student archive contains no completed assembly encoder or answer keys. Planned
light-mode figures remain explicitly planned until the separate authoring plugin
captures and verifies their real IDE states and native annotations are reviewed.
